%    Nota lateral sobre panic:
% 	panic é o último recurso do kernel: o impossível aconteceu e o
% 	kernel não sabe como proceder. No xv6, panic faz ...
\chapter{Traps e chamadas de sistema}
\label{CH:TRAP}

Existem três tipos de eventos que fazem a CPU interromper
a execução normal de instruções e transferir o controle
para um código especial que trata o evento. Uma dessas situações
é uma chamada de sistema, quando um programa de usuário
executa a instrução {\tt ecall} para solicitar que o kernel
realize alguma ação. Outra situação é uma \indextext{exceção}:
uma instrução (de usuário ou kernel) faz algo ilegal, como uma divisão
por zero ou o uso de um endereço virtual inválido. A terceira situação é uma
\indextext{interrupção} de dispositivo, quando um dispositivo sinaliza que precisa
de atenção, por exemplo, quando o hardware de disco finaliza uma requisição
de leitura ou escrita.

Este livro usa \indextext{trap} como um termo genérico para essas
situações. Tipicamente, o código que estava sendo executado no momento da
trap precisará ser retomado depois, e não deve perceber que algo
especial aconteceu. Ou seja, geralmente queremos que as traps sejam
transparentes; isso é particularmente importante para interrupções de dispositivo,
as quais o código interrompido normalmente não espera. A sequência usual é que
uma trap transfere o controle para o kernel; o kernel salva os registradores
e outros estados para que a execução possa ser retomada; o kernel
executa o código apropriado de tratamento (por exemplo, uma implementação
de chamada de sistema ou um driver de dispositivo); o kernel restaura o estado salvo e retorna
da trap; e o código original retoma do ponto em que parou.

O xv6 trata todas as traps no kernel; traps não são entregues ao código
de usuário. Tratar traps no kernel é natural no caso de chamadas de sistema.
Faz sentido para interrupções porque o isolamento exige que apenas o
kernel possa acessar dispositivos,
e porque o kernel é um mecanismo conveniente para
compartilhar dispositivos entre múltiplos processos.
Também faz sentido para exceções,
já que o xv6 responde a todas as exceções vindas do espaço do usuário
matando o programa responsável.

O tratamento de traps no xv6 ocorre em quatro estágios: ações de hardware tomadas pela
CPU RISC-V, algumas instruções em assembly que preparam o caminho para
o código em C do kernel, uma função em C que decide o que fazer com a trap,
e a rotina de serviço da chamada de sistema ou do driver de dispositivo. Embora
a semelhança entre os três tipos de trap sugira que o kernel poderia
tratá-los todos com um único caminho de código, na prática é
conveniente ter código separado para
dois casos distintos: traps vindas do espaço do usuário, e traps do espaço do kernel.
O código do kernel (em assembly ou C) que
processa uma trap é frequentemente chamado de \indextext{handler};
as primeiras instruções do handler são geralmente escritas em assembly
(em vez de C) e às vezes são chamadas de \indextext{vetor}.


\section{Mecanismo de trap do RISC-V}

Cada CPU RISC-V possui um conjunto de registradores de controle que o kernel escreve
para informar à CPU como lidar com traps, e que o kernel pode ler
para descobrir sobre uma trap que ocorreu. Os documentos do RISC-V
contêm todos os detalhes~\cite{riscv:priv}. O arquivo {\tt riscv.h}
\lineref{kernel/riscv.h:1} contém as definições utilizadas pelo xv6. Aqui está
um resumo dos registradores mais importantes:

\begin{itemize}

\item \indexcode{stvec}: O kernel escreve aqui o endereço de seu handler de trap;
  o RISC-V salta para o endereço em {\tt stvec} para tratar uma trap.

\item \indexcode{sepc}: Quando uma trap ocorre, o RISC-V salva o contador de programa
  aqui (pois o {\tt pc} será sobrescrito com o
  valor em {\tt stvec}). A instrução
  {\tt sret} (retorno de trap) copia {\tt sepc} para o
  {\tt pc}. O kernel pode escrever em {\tt sepc} para controlar para onde {\tt
    sret} retornará.

\item \indexcode{scause}: O RISC-V coloca aqui um número que descreve
a razão da trap.

\item \indexcode{sscratch}: O código do handler de trap usa {\tt sscratch}
  para evitar sobrescrever registradores de usuário antes de salvá-los.

\item \indexcode{sstatus}: O bit SIE em {\tt sstatus}
  controla se interrupções de dispositivo
  estão habilitadas. Se o kernel limpar o bit SIE, o RISC-V adiará
  interrupções de dispositivo até que o kernel o reative. O bit SPP
  indica se a trap veio do modo usuário ou supervisor,
  e controla para qual modo {\tt sret} retorna.

\end{itemize}

Os registradores acima se referem a traps tratadas no modo supervisor, e eles
não podem ser lidos ou escritos em modo usuário.

Cada CPU em um chip multi-core possui seu próprio conjunto desses registradores,
e mais de uma CPU pode estar lidando com uma trap ao mesmo tempo.

Quando precisa forçar uma trap, o hardware RISC-V faz o
seguinte para todos os tipos de trap:

\begin{enumerate}

\item Se a trap for uma interrupção de dispositivo e o bit SIE de {\tt sstatus}
  estiver limpo, não faça mais nada.

\item Desabilita interrupções limpando o bit SIE em {\tt sstatus}.

\item Copia o {\tt pc} para {\tt sepc}.

\item Salva o modo atual (usuário ou supervisor) no bit SPP de {\tt sstatus}.

\item Define {\tt scause} para refletir a causa da trap.

\item Define o modo para supervisor.

\item Copia {\tt stvec} para o {\tt pc}.

\item Começa a execução no novo {\tt pc}.

\end{enumerate}

Note que a CPU não muda para a tabela de páginas do kernel, não
muda para uma pilha no kernel, e não salva registradores além
do {\tt pc}. O software do kernel deve realizar essas tarefas.
Uma das razões para que a CPU faça o mínimo possível durante uma trap
é oferecer flexibilidade ao software; por exemplo, alguns sistemas operacionais
evitam mudar a tabela de páginas em certas situações para melhorar
o desempenho da trap.

Vale a pena pensar se algum dos passos listados acima poderia
ser omitido, talvez em busca de traps mais rápidas. Embora existam
situações em que uma sequência mais simples funcione, muitos dos passos
seriam perigosos de omitir em geral. Por exemplo, suponha que a
CPU não mudasse o contador de programa. Então uma trap vinda do espaço do usuário poderia
mudar para o modo supervisor ainda executando instruções de usuário. Essas
instruções poderiam violar o isolamento entre usuário e kernel, por exemplo
modificando o registrador {\tt satp} para apontar para uma tabela de páginas que
permite acesso a toda a memória física. Por isso, é importante que
a CPU altere o endereço de instrução para um especificado pelo kernel, ou seja, {\tt
  stvec}.

\section{Traps do espaço do usuário}

O Xv6 lida com traps de maneira diferente, dependendo de
se a trap ocorre enquanto
executa no núcleo ou no código do usuário. Aqui está a
história das traps do código do usuário; a Seção~\ref{s:ktraps}
descreve traps do código do núcleo.

Uma trap pode ocorrer enquanto executa no espaço do usuário se o
programa do usuário fizer uma
chamada de sistema (instrução {\tt ecall}), ou fizer algo
ilegal, ou se um dispositivo interromper.
O caminho de alto nível de uma trap do espaço do usuário é
{\tt uservec}
\lineref{kernel/trampoline.S:/^uservec/},
depois {\tt usertrap}
\lineref{kernel/trap.c:/^usertrap/};
e ao retornar,
{\tt usertrapret}
\lineref{kernel/trap.c:/^usertrapret/}
e então
{\tt userret}
\lineref{kernel/trampoline.S:/^userret/}.

% falar sobre por que RISC-V não troca tabelas de páginas

Uma grande restrição no design do tratamento de traps do xv6 é o fato
de que o hardware RISC-V não troca tabelas de páginas quando força uma
trap. Isso significa que o endereço do manipulador de trap
em {\tt stvec} deve ter um mapeamento válido na tabela de páginas do usuário,
já que essa é a tabela de páginas em vigor
quando o código de tratamento de trap começa a ser executado. Além disso, o código de
tratamento de trap do xv6 precisa trocar para a tabela de páginas do núcleo; para poder
continuar executando após essa troca, a tabela de páginas do núcleo
também deve ter um mapeamento para o manipulador apontado por {\tt stvec}.

O Xv6 satisfaz esses requisitos usando uma página de \indextext{trampolim}.
A página de trampolim contém {\tt uservec}, o código de tratamento de trap do xv6
que {\tt stvec} aponta. A página de trampolim está mapeada na
tabela de páginas de cada processo no endereço
\indexcode{TRAMPOLINE},
que está no topo do espaço de endereços virtual para que fique
acima da memória que os programas usam para si mesmos.
A página de trampolim também está mapeada no endereço {\tt TRAMPOLINE}
na tabela de páginas do núcleo. Veja a Figura~\ref{fig:as} e
a Figura~\ref{fig:xv6_layout}. Como a página de trampolim está mapeada na
tabela de páginas do usuário, as traps podem começar
a ser executadas lá em modo supervisor. Como a página de trampolim está
mapeada no mesmo endereço no espaço de endereços do núcleo, o manipulador de trap
pode continuar a executar após trocar para a tabela de páginas do núcleo.

O código do manipulador de trap {\tt uservec} está em {\tt trampoline.S}
\lineref{kernel/trampoline.S:/^uservec/}.
Quando {\tt uservec} começa, todos os 32 registradores contêm valores pertencentes ao
código do usuário interrompido. Esses 32 valores precisam ser salvos em algum lugar
na memória, para que mais tarde
o núcleo possa restaurá-los antes de retornar ao espaço do usuário.
Armazenar na memória requer o uso de um registrador
para manter o endereço,
mas neste ponto não há registradores de uso geral disponíveis!
Felizmente, o RISC-V fornece uma ajuda na
forma do registrador {\tt sscratch}. A instrução {\tt csrw} no
início de {\tt uservec} salva {\tt a0} em {\tt
  sscratch}. Agora 
{\tt uservec} tem
um registrador ({\tt a0}) para trabalhar.

A próxima tarefa de {\tt uservec} é salvar os 32 registradores do usuário.
O núcleo aloca, para cada processo, uma página de memória para uma
estrutura {\tt trapframe} que (entre outras coisas) tem espaço para
salvar os 32 registradores do usuário
\lineref{kernel/proc.h:/^struct.trapframe/}. Como {\tt satp} ainda
se refere à tabela de páginas do usuário, {\tt uservec} precisa que o trapframe seja
mapeado no espaço de endereços do usuário. O Xv6 mapeia o trapframe de cada processo
no endereço virtual {\tt TRAPFRAME} na tabela de páginas do usuário desse processo;
{\tt TRAPFRAME} está
logo abaixo de {\tt TRAMPOLINE}. O {\tt p->trapframe} do processo também
aponta para o
trapframe, embora em seu endereço físico para que o núcleo possa usá-lo
através da tabela de páginas do núcleo.

Assim, {\tt uservec} carrega o endereço {\tt TRAPFRAME} em {\tt a0}
e salva todos os registradores do usuário lá,
incluindo o {\tt a0} do usuário, lido de volta de {\tt sscratch}.

O {\tt trapframe} contém o endereço da
pilha do núcleo do processo atual, o hartid da CPU atual, o endereço da função
{\tt usertrap},
e o endereço da tabela de páginas do núcleo. {\tt uservec}
recupera esses valores, troca {\tt satp} para a tabela de páginas do núcleo,
e salta para {\tt usertrap}.

A tarefa de {\tt usertrap} é determinar
a causa da trap, processá-la e retornar
\lineref{kernel/trap.c:/^usertrap/}.
Ele primeiro altera {\tt stvec} para que
uma trap enquanto no núcleo seja tratada por
{\tt kernelvec} em vez de {\tt uservec}.
Ele salva o registrador {\tt sepc} (o contador de programa do usuário salvo),
porque 
{\tt usertrap} pode chamar \lstinline{yield} para
trocar para o thread do núcleo de outro processo, e esse processo pode retornar
ao espaço do usuário, durante o qual ele modificará \lstinline{sepc}.
Se a trap for uma chamada de sistema, {\tt usertrap} chama {\tt syscall} para
tratá-la;
se for uma interrupção de dispositivo, {\tt devintr};
caso contrário, é uma exceção, e o núcleo finaliza o
processo com falha.
O caminho da chamada de sistema adiciona quatro ao contador de programa do usuário salvo
porque o RISC-V, no caso de uma chamada de sistema,
deixa o ponteiro do programa apontando para a instrução {\tt ecall}
mas o código do usuário precisa retomar a execução na instrução subsequente.
No caminho de saída, {\tt usertrap} verifica se o processo foi
finalizado ou se deve ceder a CPU (se essa trap for uma interrupção de temporizador).

O primeiro passo para retornar ao espaço do usuário é a chamada para {\tt usertrapret}
\lineref{kernel/trap.c:/^usertrapret/}.
Essa função configura os registradores de controle do RISC-V para se preparar para uma
futura trap do espaço do usuário: configurando {\tt stvec}
para {\tt uservec} e preparando os campos do trapframe que
{\tt uservec} depende.
{\tt usertrapret} define {\tt sepc} para o contador de programa do usuário previamente
salvo.
No final, {\tt usertrapret}
chama {\tt userret} na página de trampolim que está mapeada em
tanto na tabela de páginas do usuário quanto na do núcleo; a razão é que o código assembly
em {\tt userret} trocará tabelas de páginas.

A chamada de {\tt usertrapret} para {\tt userret} passa 
um ponteiro para a tabela de páginas do usuário do processo em {\tt a0}
\lineref{kernel/trampoline.S:/^userret/}.
{\tt userret} troca {\tt satp} para a tabela de páginas do usuário do processo.
Lembre-se de que a tabela de páginas do usuário mapeia tanto a página de trampolim
quanto {\tt TRAPFRAME}, mas nada mais do núcleo.
O mapeamento da página de trampolim no mesmo
endereço virtual nas tabelas de páginas do usuário e do núcleo permite
que {\tt userret} continue executando após mudar {\tt satp}.
A partir deste ponto, os únicos dados que {\tt userret} pode usar são
os conteúdos dos registradores e o conteúdo do trapframe.
{\tt userret} carrega o endereço {\tt TRAPFRAME} em {\tt a0},
restaura os registradores do usuário salvos do trapframe via {\tt a0},
restaura o {\tt a0} do usuário salvo,
e executa {\tt sret} para retornar ao espaço do usuário.

\section{Código: Chamando chamadas de sistema}

O Capítulo~\ref{CH:FIRST} terminou com 
\indexcode{initcode.S}
invocando a chamada de sistema {\tt exec}
\lineref{user/initcode.S:/SYS_exec/}.
Vamos analisar como a chamada do usuário
chega à implementação da chamada de sistema {\tt exec} no núcleo.

{\tt initcode.S}
coloca os argumentos para
\indexcode{exec}
nos registradores {\tt a0} e {\tt a1}, e coloca o
número da chamada de sistema em
\texttt{a7}.
Os números das chamadas de sistema correspondem às entradas no array {\tt syscalls},
uma tabela de ponteiros de função
\lineref{kernel/syscall.c:/syscalls/}.
A instrução \lstinline{ecall} provoca uma trap no núcleo
e faz com que
{\tt uservec},
{\tt usertrap} e, em seguida, {\tt syscall} sejam executados, como vimos acima.

\indexcode{syscall}
\lineref{kernel/syscall.c:/^syscall/} 
recupera o número da chamada de sistema do {\tt a7} salvo no trapframe
e o usa para indexar no {\tt syscalls}.
Para a primeira chamada de sistema, 
\texttt{a7}
contém 
\indexcode{SYS_exec}
\lineref{kernel/syscall.h:/SYS_exec/},
resultando em uma chamada para a função de implementação da chamada de sistema
\lstinline{sys_exec}.

Quando \lstinline{sys_exec} retorna,
\lstinline{syscall}
registra seu valor de retorno em
\lstinline{p->trapframe->a0}.
Isso fará com que a chamada original do espaço do usuário 
{\tt exec()} retorne esse valor, uma vez que a
convenção de chamada C no RISC-V coloca valores de retorno em {\tt a0}.
Chamadas de sistema geralmente retornam números negativos para indicar
erros, e zero ou números positivos para sucesso.
Se o número da chamada de sistema for inválido,
\lstinline{syscall}
imprime um erro e retorna $-1$.

\section{Código: Argumentos da chamada de sistema}

As implementações de chamadas de sistema no núcleo precisam encontrar os argumentos
passados pelo código do usuário. Como o código do usuário chama funções wrapper de chamadas de sistema,
os argumentos estão inicialmente onde a convenção de chamada C do RISC-V os coloca: nos registradores.
O código de trap do núcleo salva os registradores do usuário no trap frame do processo atual, onde o código do núcleo pode encontrá-los.
As funções do núcleo
\lstinline{argint},
\lstinline{argaddr},
e
\lstinline{argfd}
recuperam o 
\textit{n} 'ésimo 
argumento da chamada de sistema
do trap frame
como um inteiro, ponteiro ou descritor de arquivo.
Todos eles chamam {\tt argraw} para recuperar o registrador de usuário salvo apropriado
\lineref{kernel/syscall.c:/^argraw/}.

Algumas chamadas de sistema passam ponteiros como argumentos, e o núcleo deve usar
esses ponteiros para ler ou escrever na memória do usuário. A chamada de sistema {\tt exec}, por exemplo, passa ao núcleo um array de ponteiros
referindo-se a argumentos de string no espaço do usuário.
Esses ponteiros apresentam
duas dificuldades. Primeiro, o programa do usuário pode estar com defeito ou ser malicioso, e
pode passar ao núcleo um ponteiro inválido ou um ponteiro destinado a enganar
o núcleo ao acessar a memória do núcleo em vez da memória do usuário.
Em segundo lugar, os mapeamentos da tabela de páginas do núcleo xv6 não são os mesmos que os
mapeamentos da tabela de páginas do usuário, então o núcleo não pode usar instruções comuns para carregar ou armazenar a partir de endereços fornecidos pelo usuário.

O núcleo implementa funções que transferem dados de forma segura para e
a partir de endereços fornecidos pelo usuário.
{\tt fetchstr} é um exemplo \lineref{kernel/syscall.c:/^fetchstr/}.
Chamadas de sistema de arquivos, como
{\tt exec}, usam {\tt fetchstr} para recuperar argumentos de nome de arquivo em string do espaço do usuário.
\lstinline{fetchstr} chama \lstinline{copyinstr}
para fazer o trabalho pesado.

\indexcode{copyinstr}
\lineref{kernel/vm.c:/^copyinstr/} copia até \lstinline{max} bytes para
\lstinline{dst} a partir do endereço virtual \lstinline{srcva} na tabela de páginas do usuário \lstinline{pagetable}.
Como \lstinline{pagetable} {\it não} é a tabela de páginas atual,
\lstinline{copyinstr} usa {\tt walkaddr}
(que chama {\tt walk}) para procurar
\lstinline{srcva} em
\lstinline{pagetable}, resultando no
endereço físico \lstinline{pa0}.
A tabela de páginas do núcleo mapeia toda a RAM física 
em endereços virtuais que são iguais ao endereço físico da RAM.
Isso permite que
{\tt copyinstr} copie diretamente os bytes da string de {\tt pa0} para {\tt dst}.
{\tt walkaddr} 
\lineref{kernel/vm.c:/^walkaddr/}
verifica se o endereço virtual fornecido pelo usuário é parte
do espaço de endereços do usuário do processo, para que os programas
não possam enganar o núcleo a ler outra memória.
Uma função similar, {\tt copyout}, copia dados do
núcleo para um endereço fornecido pelo usuário.

\section{Traps do espaço do núcleo}
\label{s:ktraps}

O Xv6 lida com traps do código do núcleo de maneira diferente
do que com traps do código do usuário.
Ao entrar no núcleo, {\tt usertrap} aponta {\tt stvec}
para o código assembly em {\tt kernelvec}
\lineref{kernel/kernelvec.S:/^kernelvec/}.
Como {\tt kernelvec} só é executado se
xv6 já estiver no núcleo, {\tt kernelvec} pode confiar
que {\tt satp} está definido para a tabela de páginas do núcleo, e que o
ponteiro da pilha refere-se a uma pilha de núcleo válida.
{\tt kernelvec} empilha todos os 32 registradores na pilha,
dos quais ele os restaurará mais tarde
para que o código do núcleo interrompido possa retomar sem interrupções.

{\tt kernelvec} salva os registradores na pilha do thread do núcleo interrompido,
o que faz sentido porque os valores dos registradores pertencem a
esse thread. Isso é particularmente importante se a trap causar uma
troca para um thread diferente -- nesse caso, a trap realmente
retornará da pilha do novo thread, deixando os registradores salvos do
thread interrompido seguros em sua pilha.

{\tt kernelvec} salta para {\tt kerneltrap}
\lineref{kernel/trap.c:/^kerneltrap/} após salvar os registradores.
{\tt kerneltrap} está preparado para dois tipos de traps:
interrupções de dispositivo e exceções. Ele chama
{\tt devintr}
\lineref{kernel/trap.c:/^devintr/}
para verificar e lidar com a primeira.
Se a trap não for uma interrupção de dispositivo, deve ser uma exceção,
e isso é sempre um erro fatal se ocorrer no núcleo do xv6;
o núcleo chama \lstinline{panic} e para de executar.

Se {\tt kerneltrap} foi chamado devido a uma interrupção de temporizador, e um
thread do núcleo de um processo está em execução (em oposição a um thread de escalonador),
{\tt kerneltrap} chama {\tt yield} para dar a outros threads a chance de
executar. Em algum momento, um desses threads irá ceder, permitindo que nosso thread
e seu {\tt kerneltrap} sejam retomados novamente.
O Capítulo~\ref{CH:SCHED} explica o que acontece em {\tt yield}.

Quando o trabalho de {\tt kerneltrap} está concluído, ele precisa retornar ao código
que foi interrompido pela trap. Como um {\tt yield} pode ter
perturbado {\tt sepc} e o modo anterior em {\tt sstatus},
{\tt kerneltrap} os salva quando começa. Ele agora restaura esses
registradores de controle e retorna para {\tt kernelvec}
\lineref{kernel/kernelvec.S:/call.kerneltrap$/}.
{\tt kernelvec} desempilha os registradores salvos da pilha e
executa {\tt sret}, que copia {\tt sepc} para {\tt pc}
e retoma o código do núcleo interrompido.

Vale a pena pensar em como o retorno da trap acontece se
{\tt kerneltrap} chamou {\tt yield} devido a uma interrupção de temporizador.

O Xv6 define o {\tt stvec} de uma CPU para {\tt kernelvec} quando essa CPU
entra no núcleo a partir do espaço do usuário; você pode ver isso em {\tt usertrap}
\lineref{kernel/trap.c:/stvec.*kernelvec/}.
Há uma janela de tempo em que o núcleo começou a ser executado
mas {\tt stvec} ainda está definido como {\tt uservec}, e é crucial que 
nenhuma interrupção de dispositivo ocorra durante essa janela.
Felizmente, o RISC-V sempre desativa interrupções quando começa
a lidar com uma trap, e {\tt usertrap} não as reativa até
depois de definir {\tt stvec}.

\section{Exceções de falha de página}
\label{sec:pagefaults}

A resposta do Xv6 a exceções é bastante simples: se uma exceção acontece
no espaço do usuário, o núcleo finaliza o processo com falha. Se uma
exceção acontece no núcleo, o núcleo entra em pânico. Sistemas operacionais
reais frequentemente respondem de maneiras muito mais interessantes.

Como exemplo,
muitos núcleos usam falhas de página para implementar
\indextext{fork de cópia sob demanda (COW)}.
Para explicar o fork de cópia sob demanda, considere o \lstinline{fork} do xv6,
descrito no Capítulo~\ref{CH:MEM}.
O \lstinline{fork} faz com que o conteúdo inicial da memória do filho
seja o mesmo que o do pai no momento do fork.
O Xv6 implementa o fork com
\lstinline{uvmcopy}
\lineref{kernel/vm.c:/^uvmcopy/},
que aloca memória física para
o filho e copia a memória do pai para ele.
Seria mais eficiente se o filho e o pai pudessem compartilhar 
a memória física do pai.
Uma implementação direta disso não funcionaria, no entanto,
já que causaria a interrupção da execução um do outro com suas gravações na pilha e heap compartilhados.

Pai e filho podem compartilhar com segurança a memória física por meio
do uso apropriado de permissões de tabela de páginas e falhas de página.
A CPU gera uma
\indextext{exceção de falha de página}
quando um endereço virtual é utilizado que não tem mapeamento
na tabela de páginas, ou tem um mapeamento cujo flag \lstinline{PTE_V}
está apagado, ou um mapeamento cujos bits de permissão
(\lstinline{PTE_R},
\lstinline{PTE_W},
\lstinline{PTE_X},
\lstinline{PTE_U})
proíbem a operação sendo tentada.
O RISC-V distingue três
tipos de falha de página: falhas de página de carregamento (causadas por instruções de carregamento),
falhas de página de armazenamento (causadas por instruções de armazenamento),
e falhas de página de instrução (causadas por buscas de instruções a serem executadas). O
registrador \lstinline{scause} indica o tipo da
falha de página e o registrador \indexcode{stval} contém o endereço
que não pôde ser traduzido.

O plano básico no fork COW é que o pai e o filho compartilhem inicialmente
todas as páginas físicas, mas cada um as mapeie como somente leitura (com o
flag \lstinline{PTE_W} apagado).
Pai e filho podem ler da memória física compartilhada.
Se um dos dois gravar em uma determinada página,
a CPU RISC-V gera uma exceção de falha de página.
O manipulador de trap do núcleo responde alocando uma
nova página de memória física e copiando para ela
a página física que o endereço falhado mapeia.
O núcleo altera o PTE relevante na tabela de páginas do processo com falha
para apontar para a cópia e permitir gravações, além de leituras,
e então retoma o processo com falha na instrução que causou a falha. Como o 
PTE agora permite gravações, a instrução reexecutada
será executada sem uma falha. A cópia sob demanda requer controle
para ajudar a decidir quando as páginas físicas podem ser liberadas, já que cada página pode
ser referenciada por um número variável de tabelas de páginas dependendo da história de
forks, falhas de página, execuções e saídas. Esse controle permite
uma otimização importante: se um processo sofrer uma falha de página de armazenamento
e a página física for referida apenas pela tabela de páginas daquele processo,
nenhuma cópia é necessária.

A cópia sob demanda torna o \lstinline{fork} mais rápido, já que o \lstinline{fork}
não precisa copiar memória. Parte da memória terá que ser copiada
mais tarde, quando for gravada, mas é comum que a maior parte da
memória nunca precise ser copiada.
Um exemplo comum é
\lstinline{fork} seguido de \lstinline{exec}:
poucas páginas podem ser escritas após o \lstinline{fork},
mas então o \lstinline{exec} do filho libera
a maior parte da memória herdada do pai.
A cópia sob demanda \lstinline{fork} elimina a necessidade de
copiar essa memória.
Além disso, o fork COW é transparente:
nenhuma modificação nas aplicações é necessária para
que elas se beneficiem.

A combinação de tabelas de páginas e falhas de página abre uma ampla gama
de possibilidades interessantes além do fork COW. Outra
característica amplamente utilizada é chamada de \indextext{alocação preguiçosa}, que tem
duas partes. Primeiro, quando uma aplicação solicita mais memória chamando
\lstinline{sbrk}, o núcleo anota o aumento de tamanho, mas não
aloca memória física e não cria PTEs para o novo intervalo de
endereços virtuais. Em segundo lugar, em uma falha de página em um dos novos
endereços, o núcleo aloca uma página de memória física e a mapeia
na tabela de páginas.
Assim como o fork COW,
o núcleo pode implementar a alocação preguiçosa de forma transparente para as aplicações.

Como as aplicações costumam solicitar mais memória do que realmente precisam, a alocação
preguiçosa é vantajosa: o núcleo não precisa fazer nenhum trabalho para
páginas que a aplicação nunca usa.
Além disso, se a aplicação está
pedindo para aumentar o espaço de endereços significativamente, então \lstinline{sbrk} 
sem alocação preguiçosa é
caro: se uma aplicação pede um gigabyte de memória,
o núcleo precisa alocar e zerar 262.144 páginas de 4096 bytes.
A alocação preguiçosa permite que esse custo
seja diluído ao longo do tempo. Por outro lado, a alocação preguiçosa
incorre na sobrecarga extra de falhas de página, que envolvem uma transição entre usuário/núcleo.
Os sistemas operacionais podem reduzir esse custo alocando um lote de 
páginas consecutivas por
falha de página, em vez de uma única página, e especializando o código de entrada/saída do núcleo
para tais falhas de página.

Outra característica amplamente utilizada que explora falhas de página é
\indextext{paginação sob demanda}. No \lstinline{exec}, o xv6 carrega todo 
o texto
e dados de uma aplicação na memória antes de iniciar
a aplicação. Como as aplicações
podem ser grandes e ler do disco leva tempo, esse custo de inicialização pode
ser notável para os usuários. Para
diminuir o tempo de inicialização, um núcleo moderno não carrega inicialmente
o arquivo executável na memória, mas apenas cria a tabela de páginas do usuário com
todos os PTEs marcados como inválidos. O núcleo inicia o programa;
cada vez que o programa usa uma página pela primeira vez, ocorre uma falha de página,
e em resposta
o núcleo lê o conteúdo da página do disco e
mapeia na memória do espaço de endereços do usuário. Assim como o fork COW e a alocação
preguiçosa, o núcleo pode implementar essa característica de forma transparente para
as aplicações.

Os programas em execução em um computador podem precisar de mais memória do que a
RAM disponível. Para lidar com isso de forma adequada, o sistema operacional pode
implementar \indextext{paginação para disco}. A ideia é armazenar apenas uma
fração das páginas do usuário na RAM e armazenar o restante no disco em uma
\indextext{área de paginação}. O núcleo marca os PTEs que correspondem à
memória armazenada na área de paginação (e, portanto, não na RAM) como inválidos. Se
uma aplicação tentar usar uma das páginas que foi {\it paginada
  para fora} para o disco, a aplicação sofrerá uma falha de página, e a página
precisará ser {\it paginada para dentro}: o manipulador de trap do núcleo alocará uma página
de RAM física, lerá a página do disco para a RAM e modificará o
PTE relevante para apontar para a RAM.

O que acontece se uma página precisar ser paginada para dentro, mas não houver RAM física
livre? Nesse caso, o núcleo deve primeiro liberar uma página física
paginando-a para fora ou {\it expulsando-a} para a área de paginação no disco, e
marcando os PTEs que se referem a essa página física como inválidos. A expulsão
é cara, então a paginação tem melhor desempenho se for infrequente: se
as aplicações usarem apenas um subconjunto de suas páginas de memória e a união dos
subconjuntos couber na RAM. Essa propriedade é frequentemente referida como tendo
boa localidade de referência. Assim como muitas técnicas de memória virtual,
os núcleos geralmente implementam a paginação para disco de uma maneira que é transparente
para as aplicações.

Os computadores frequentemente operam com pouca ou nenhuma memória física {\it livre},
independentemente de quanta RAM o hardware fornece. Por exemplo, provedores de nuvem
multiplexam muitos clientes em uma única máquina para utilizar seus
custos de hardware de forma eficaz. Como outro exemplo, usuários executam muitos
aplicativos em smartphones com uma pequena quantidade de memória física. Em
tais configurações, alocar uma página pode exigir primeiro expulsar uma página existente. Assim, quando a memória física livre é escassa, a alocação é
cara.

A alocação preguiçosa e a paginação sob demanda são particularmente vantajosas quando
a memória livre é escassa e os programas utilizam ativamente apenas uma fração de
sua memória alocada. Essas técnicas também podem evitar o trabalho
desperdiçado quando uma página é alocada ou carregada, mas nunca usada ou
expulsa antes que possa ser utilizada.

Outras características que combinam paginação e exceções de falha de página incluem
a extensão automática de pilhas e \indextext{arquivos mapeados na memória},
que são arquivos que um programa mapeou em seu espaço de endereços usando
a chamada de sistema \texttt{mmap} para que o programa possa lê-los e escrevê-los
usando instruções de carregamento e armazenamento.

\section{Mundo real}

O trampolim e o trapframe podem parecer excessivamente complexos. Uma força motriz é que o RISC-V intencionalmente faz o mínimo possível ao forçar uma trap, para permitir a possibilidade de um 
tratamento de trap muito rápido, o que se mostra importante. Como resultado, as primeiras instruções do manipulador de trap do núcleo efetivamente têm que ser executadas no ambiente do usuário: a 
tabela de páginas do usuário e o conteúdo dos registradores do usuário.
E o manipulador de trap inicialmente ignora fatos úteis, como a
identidade do processo que está em execução ou o endereço da tabela de páginas do núcleo. Uma solução é possível porque o RISC-V fornece lugares protegidos nos quais o núcleo pode armazenar 
informações antes de entrar no espaço do usuário: o registrador {\tt sscratch}, e entradas de tabela de páginas do usuário
que apontam para a memória do núcleo, mas são protegidas pela falta de \lstinline{PTE_U}.
O trampolim e o trapframe do Xv6 exploram esses recursos do RISC-V.

A necessidade de páginas de trampolim especiais poderia ser eliminada se a memória do núcleo
fosse mapeada na tabela de páginas do usuário de cada processo (com
\lstinline{PTE_U} apagado).
Isso também eliminaria a necessidade de uma troca de tabela de páginas ao fazer uma trap do
espaço do usuário para o núcleo. Isso, por sua vez, permitiria que implementações de chamadas de sistema no núcleo aproveitassem a memória do usuário do processo atual sendo mapeada, permitindo que 
o código do núcleo desreferencie diretamente ponteiros do usuário. Muitos sistemas operacionais usaram essas ideias para
aumentar a eficiência. O Xv6 as evita para reduzir as chances de
erros de segurança no núcleo devido ao uso inadvertido de ponteiros do usuário,
e para reduzir alguma complexidade que seria necessária para garantir que
endereços virtuais do usuário e do núcleo não se sobreponham.

Sistemas operacionais de produção implementam fork de cópia sob demanda, alocação
preguiçosa, paginação sob demanda, paginação para disco, arquivos mapeados na memória, etc.
Além disso, sistemas operacionais de produção tentam armazenar
algo útil em todas as áreas da memória física, normalmente armazenando em cache
o conteúdo de arquivos na memória que não está sendo usado por processos.

Sistemas operacionais de produção também fornecem às aplicações chamadas de sistema
para gerenciar seus espaços de endereços e implementar seu próprio
tratamento de falhas de página através das chamadas de sistema {\tt mmap}, {\tt munmap} e {\tt
  sigaction}, além de fornecer chamadas para fixar a memória
na RAM (veja {\tt mlock}) e para aconselhar o núcleo sobre como uma aplicação
planeja usar sua memória (veja {\tt madvise}).

\section{Exercícios}

\begin{enumerate}

\item As funções {\tt copyin} e {\tt copyinstr} percorrem a tabela de páginas do usuário em software. Configure a tabela de páginas do núcleo para que o núcleo tenha o programa do usuário mapeado, 
e {\tt copyin} e {\tt
    copyinstr} possam usar {\tt memcpy} para copiar argumentos de chamadas de sistema para o espaço do núcleo, confiando no hardware para percorrer a tabela de páginas.

\item Implemente alocação de memória preguiçosa.

\item Implemente o fork COW.

\item Existe uma maneira de eliminar o mapeamento especial da página {\tt
  TRAPFRAME} em cada espaço de endereços do usuário? Por
  exemplo, poderia
  {\tt uservec} ser modificado para simplesmente empurrar os 32 registradores do usuário
  na pilha do núcleo, ou armazená-los na estrutura {\tt proc}?

\item O xv6 poderia ser modificado para eliminar o mapeamento especial da página {\tt
  TRAMPOLINE}?

\item Implemente {\tt mmap}.

\end{enumerate}
